\section{Analisis Permasalahan}

Analisis kesenjangan DecMed terhadap PGHD
Tujuan: mengikat langsung ke Bab 2 (DecMed \& PGHD).

Isi yang bisa kamu bahas:

\begin{enumerate}
    \item Ringkas kembali kemampuan DecMed saat ini (kontrol akses, PRE, IPFS, IOTA).
    
    \item Tunjukkan secara eksplisit bahwa:
    \begin{itemize}
        \item pipeline data sekarang provider-centric; hanya episodik dari fasyankes,
        \item belum ada aliran data: pasien → smartphone/wearable → DecMed secara kontinu.
    \end{itemize}
    
    \item Kaitkan dengan literatur PGHD–EHR: studi integrasi PGHD - EHR menunjukkan kebutuhan pipeline yang kontinu \& terstruktur.
    
    \item Output: daftar gap, misalnya:
    
    \begin{itemize}
        \item tidak ada modul ingest PGHD,
        
        \item tidak ada mekanisme offline-first di sisi pasien,
        
        \item tidak ada skema provenance khusus PGHD.
    \end{itemize}
\end{enumerate}

Ini jadi “problem statement teknis” yang Bab 3-mu jawab.